\section{问题分析}
本题包含三个不同尺度：任务在离散区域和整数开工时刻上选择，分钟级持续时间决定逐小时重叠，能源与储能则按小时结算。直接把任务汇总为平均负荷会掩盖瞬时 GPU 和功率冲突，因此本文始终保留任务级轨迹，并由轨迹重建区域--小时资源数组。

问题一的关键是防止未来信息泄漏。所有滞后特征只使用预测时点之前的数据，验证集只负责模型选择，测试集不参与任何超参数决策。问题二的组合空间远大于可用计算预算，故采用确定性增量评价：一次移动只更新受影响的区域--小时数组，但每个命名方案结束后仍对全部任务重新审计。问题三在给定设施负荷后按区域分解，以有限 SOC 状态动态规划产生不同目标偏好的储能轨迹；no-action 必须进入候选集。问题四不是在固定三套任务方案中挑选，而是对每个情景从共同基线独立重启，交替执行储能优化和由当前储能边际信号驱动的任务邻域重建。
